% !TEX options=--shell-escape

\documentclass{article}

\usepackage{multicol} % Enables multicol env.
\usepackage{graphicx} % Enables including graphics.
\usepackage[export]{adjustbox} % Allows graphics to have a `center` property.
\usepackage[pdf]{graphviz}
\usepackage{mdframed}
\usepackage{parskip}
\usepackage{amssymb}
\usepackage{titlesec} % Enables custom title sections.
\usepackage{float}
\usepackage{listings}



\newmdtheoremenv{theo}{Theorem}
\newmdtheoremenv{fact}{Fact}
\newmdtheoremenv{definition}{Definition}
\newmdtheoremenv{lemma}{Lemma}

\renewcommand\thesection{Question \arabic{section}}
\titleformat{\subsection}
{\normalfont\normalsize}{\thesubsection}{1em}{}

\renewcommand\thesubsection{(\alph{subsection})}
\titleformat{\subsection}
{\normalfont\normalsize}{\thesubsection}{1em}{}

\renewcommand\thesubsubsection{\hspace{10px}\roman{subsubsection}.}
\titleformat{\subsubsection}
{\normalfont\normalsize}{\thesubsubsection}{1em}{}

\newcommand{\YearPath}{../../../LatexConfig} % Get year level path (i.e. Senior)
\newcommand{\SemesterPath}{../../LatexConfig} % Get semester level path (i.e. Semester 1).
\newcommand{\ClassPath}{../LatexConfig} % Get class level path (i.e. CS101).
\newcommand{\AssignmentTitle}{Homework 4}
\newcommand{\AssignmentSub}{~ ~}

\begin{document}

\newcommand{\Professor}{Dr. Thorne Wolf}
\newcommand{\Course}{CS 101.1}
\newcommand{\Professor}{Dr. Thorne Wolf}
\newcommand{\Course}{CS 101.1}
\begin{titlepage}
  \begin{center}
      \vspace*{1cm}

      \Large
      \textbf{\AssignmentTitle}

      \vspace{0.5cm}
      \large
      \AssignmentSub \\
      ~\\
      \normalsize \today

      \vspace{1.5cm}

      \large
      By: \\
      ~\\
      \normalsize
      \vspace{1ex}
      \ifdefined\StudentNames \StudentNames \else Thorne Wolfenbarger\fi
      \vfill

      \vspace{0.8cm}

      \large
      Submitted to: \\
      \Professor \\
      In Partial Fulfillment of the Requirements of \\
      \Course~-~\Semester\\

      \vspace{0.8cm}

      \includegraphics[width=0.2\textwidth,center]{\YearPath/univ}
      ~\\
      College of Engineering\\
      Embry-Riddle Aeronautical University\\
      Prescott, AZ\\

  \end{center}
\end{titlepage}


\section{} % 1
Define a recursive sequence \\

\subsection{Provide the full statement of at least one theorem, as it’s stated on the presenter’s slides.}
Every $C^*$-algebra looks like a closed $*$-subalgebra of $B(H)$ for some Hilbert space $H$.

\subsection{In the theorem statement above, what type of a proposition is it? Is it conditional? Does it have suppressed quantifiers? If so, rewrite the theorem statement showing these.}
It is a universal statement because the theorem states that ``Every $C^*$-algebra...'' an example of this is:

\begin{center}
  $c(x) = \{f : x \rightarrow \mathbb{C} : $f is continuous$\}$ where x is compact, Hausdorff...
\end{center}

I suppose that the theorem could be conditional because it has to be a $C^*$-algebra.

We can write this as
\begin{center}
  $\forall C^* algebra~\exists H ~s.t. ~C^* algebra~ = ~closed~*$-subalgebra$~ \subseteq B(H)$.
\end{center}

\subsection{Write down a math term/word/notation that seems familiar to you (look out for words like injective, surjective, or notation like function notation).}
Let $(A(\mathbb{D}), \alpha)$ be a dynamical system. Then there exists a Mobius transformation of the form

\begin{center}
  $\tau(w) = \lambda (\frac{w-\mu}{1-\mu\bar w})$
\end{center}
such that $\alpha_1(f)(w) = f(\tau(w))$.

Within this I understand what it means for a system to be dynamical and for a transformation between spaces to exist.

\subsection{What do you understand this word/notation to mean?}
I understand this exerpt to mean there is some transform called a Mobius transform that exists as a function that accepts an argument $w$ and uses it to convert into something that is possibly called Mobius space.

\subsection{Using a resource (wikipedia is fine), provide a formal definition of the word or symbol as it was intended in the presentation.}
Here I look into the definition of a Mobius transformation.

In geometry and complex analysis, a Möbius transformation of the complex plane is a rational function of the form

\begin{center}
  $f(z) = \frac{az+b}{cz+d}$
\end{center}
of one complex variable $z$; here the coefficients $a,b,c,d$ are complex numbers satisfying $ab-bc \neq 0$.

"A Mobius transformation is always a bijective holomorphic function from the Riemann sphere to the Riemann sphere."

"These transformations preserve angles, map every straight line to a line or circle, and map every circle to a line or circle."

All information is probided from (Wikipedia, 2020). These exerpts provide the definition of a Mobius transformation as well as an invarient within the transformation.

\subsection{Write down a math term/word/notation that was totally unfamiliar to you.}
Irrational Rotation.

\subsection{Using a resource (wikipedia is fine), provide a formal definition of the word or symbol as it was intended in the presentation.}
In the mathematical theory of dynamicval systems, an irrational rotation is a map where $\theta$ is an irrational number. Under the identification of a circle with $R/z$ or with the interval $[0,1]$ with the boundary points glued together, this map becomes a rotation of a circle by a proportion $\theta$ of a full revolution.

\subsection{Ask a question that the presentation raised for you that you would like to know more about.}
Why isn't any orbit periodic? Within the same slide, Dr. Hamidi makes a statement "No orbit is periodic." What is a non-periodic orbit?


\newpage
\section{} % 2
Create an interesting example of a relation on a set. Your set should have at least 5 elements, and your relation should be non-trivial.

Consider the set $S=\{2, 3, 6, 9, 12, 17\}$.\\
Consider the relation $R=\{(a,b) ~s.t.~ a$ shares a divisor with $b\}$

This means that the relation $R=$\\
$\{(2,6), (2,12), (3,6), (3,9), (3,12), (6,2), (6,3), (6,9), (6,12), (9,3), (9,6), (9,12),$\\
$(2,2), (3,3), (6,6), (9,9), (12,12), (17, 17)\}$

\subsection{Determine if the relation you defined on your chosen set is reflexive, symmetric, and/or transitive}
This relation is reflexive, symmetric, but not transitive. Reflexive because each element relates to itself, symmetric because $(a,b) \in R \rightarrow (b,a) \in R$, and not transitive because $(a,b) \in R \land (b,c) \in R \not \rightarrow (a,c) \in R$. An example of failing transitivity are the elements $\{(2,12), (12,9)\}$ but there is no $(2,9)$.

\subsection{Construct the relation matrix $M_R$ for the relation. Be sure to indicate which element of the set
represents which row of the matrix.}
\begin{figure}[H]
\centering
\begin{tabular}{|c|cccccc|}
\hline
Element & 2 & 3 & 6 & 9 & 12 & 17 \\
\hline
2 & 1 & 0 & 1 & 0 & 1 & 0 \\
3 & 0 & 1 & 1 & 1 & 1 & 0 \\
6 & 1 & 1 & 1 & 1 & 1 & 0 \\
9 & 0 & 1 & 1 & 1 & 1 & 0 \\
12 & 1 & 1 & 1 & 1 & 1 & 0 \\
17 & 0 & 0 & 0 & 0 & 0 & 1 \\
\hline
\end{tabular}
\end{figure}

\subsection{Construct a graph that corresponds to your relation. Be sure to indicate which element of the
set represents which vertex of the graph.}


\begin{figure}[H]
\centering
\includegraphics[width=4.5in]{RelationGraph}
\end{figure}

\subsection{Describe the relationship between your set relation, the relation matrix, the graph, and the adjacency matrix for the graph.}

The relation matrix $M_R$ is exactly the adjacency matrix for the resultant graph. A notable characteristic about the relation is that since it is symmetric we can represent the relation graph as an undirected graph. The reflexiveness of the relation is readily apparent since we can see the self-loop on every vertex within the graph. 

These three items, the relation, the relation matrix, and the resultant graph all communicate exactly the same information about this relation on the set $S$.

\newpage
\section{} % 3

\subsection{Provide the full statement of at least one theorem, as it’s stated on the presenter’s slides.}
Every $C^*$-algebra looks like a closed $*$-subalgebra of $B(H)$ for some Hilbert space $H$.

\subsection{In the theorem statement above, what type of a proposition is it? Is it conditional? Does it have suppressed quantifiers? If so, rewrite the theorem statement showing these.}
It is a universal statement because the theorem states that ``Every $C^*$-algebra...'' an example of this is:

\begin{center}
  $c(x) = \{f : x \rightarrow \mathbb{C} : $f is continuous$\}$ where x is compact, Hausdorff...
\end{center}

I suppose that the theorem could be conditional because it has to be a $C^*$-algebra.

We can write this as
\begin{center}
  $\forall C^* algebra~\exists H ~s.t. ~C^* algebra~ = ~closed~*$-subalgebra$~ \subseteq B(H)$.
\end{center}

\subsection{Write down a math term/word/notation that seems familiar to you (look out for words like injective, surjective, or notation like function notation).}
Let $(A(\mathbb{D}), \alpha)$ be a dynamical system. Then there exists a Mobius transformation of the form

\begin{center}
  $\tau(w) = \lambda (\frac{w-\mu}{1-\mu\bar w})$
\end{center}
such that $\alpha_1(f)(w) = f(\tau(w))$.

Within this I understand what it means for a system to be dynamical and for a transformation between spaces to exist.

\subsection{What do you understand this word/notation to mean?}
I understand this exerpt to mean there is some transform called a Mobius transform that exists as a function that accepts an argument $w$ and uses it to convert into something that is possibly called Mobius space.

\subsection{Using a resource (wikipedia is fine), provide a formal definition of the word or symbol as it was intended in the presentation.}
Here I look into the definition of a Mobius transformation.

In geometry and complex analysis, a Möbius transformation of the complex plane is a rational function of the form

\begin{center}
  $f(z) = \frac{az+b}{cz+d}$
\end{center}
of one complex variable $z$; here the coefficients $a,b,c,d$ are complex numbers satisfying $ab-bc \neq 0$.

"A Mobius transformation is always a bijective holomorphic function from the Riemann sphere to the Riemann sphere."

"These transformations preserve angles, map every straight line to a line or circle, and map every circle to a line or circle."

All information is probided from (Wikipedia, 2020). These exerpts provide the definition of a Mobius transformation as well as an invarient within the transformation.

\subsection{Write down a math term/word/notation that was totally unfamiliar to you.}
Irrational Rotation.

\subsection{Using a resource (wikipedia is fine), provide a formal definition of the word or symbol as it was intended in the presentation.}
In the mathematical theory of dynamicval systems, an irrational rotation is a map where $\theta$ is an irrational number. Under the identification of a circle with $R/z$ or with the interval $[0,1]$ with the boundary points glued together, this map becomes a rotation of a circle by a proportion $\theta$ of a full revolution.

\subsection{Ask a question that the presentation raised for you that you would like to know more about.}
Why isn't any orbit periodic? Within the same slide, Dr. Hamidi makes a statement "No orbit is periodic." What is a non-periodic orbit?


\newpage
\section{} % 4
Compute each of the following values, and show your work.
\subsection{($177$ mod $31\cdot 270$ mod $31$) mod 31}
\begin{center}
\begin{tabular}{cl}
$31\cdot 270 = 8370 > 177$ \\
177 mod 8370 = 177 & This is because 8370 is to large to divide 177 even once. \\
(177 mod 31) mod 31 = 22 mod 31 & The remainder after division is 22 \\
22 mod 31 = 22 & By the same reasoning as before. 31 does not divide 22
\end{tabular}
\end{center}

The answer is 22.

\subsection{($7^3$ mod 23)$^2$ mod 31}
\begin{center}
\begin{tabular}{clc}
($7^3$ mod 23)$^2$ mod 31 & & $\iff$ \\
$[(49 ~\rm{mod}~ 23)(7 ~\rm{mod}~ 23) ~\rm{mod}~ 23]^2$ mod 31 & By the modular multiplication identity & $\iff$ \\
$[(3)(7) ~\rm{mod}~ 23]^2 mod 31$ & Simplifying & $\iff$ \\
$21^2 ~\rm{mod}~ 31$ & Evaluating and simplifying & $\iff$ \\
$441 ~\rm{mod}~ 31 = 7$ & Evaluating
\end{tabular}
\end{center}

The answer is 7.

\subsection{$(99^2 ~\rm{mod}~ 32)^3 ~\rm{mod}~ 15$}

\begin{center}
\begin{tabular}{clc}
$(99^2 ~\rm{mod}~ 32)^3 ~\rm{mod}~ 15$ & & $\iff$ \\
$[(99 ~\rm{mod}~ 32)(99 ~\rm{mod}~ 32) ~\rm{mod}~ 32]^3 ~\rm{mod}~ 15$ & Modular multiplication rule & $\iff$ \\
$(3\cdot 3 ~\rm{mod}~ 32)^3 ~\rm{mod}~ 15$ & Evaluation & $\iff$ \\
$9^3 ~\rm{mod}~ 15$ & Evaluation & $\iff$ \\
$729 ~\rm{mod}~ 15 = 9$ & Evaluation
\end{tabular}
\end{center}

The answer is 9.

\subsection{Prove the following: if $a,b,m,n \in \mathbb{Z}, m \geq 1, n \geq 2$, and $a \equiv b ~\rm{mod}~ n$, then $a^m \equiv b^m ~\rm{mod}~ n$.}

\begin{center}
\begin{tabular}{p{0.6\linewidth}p{0.4\linewidth}c}
    $a,b,m,n \in \mathbb{Z}$ & Given as the domain for all used variables \\
    $m \geq 1, n \geq 2$ & Given \\
    $a \equiv b ~\rm{mod}~ n$ & Given \\
    $a^m = a \cdot a \cdot a\ldots \cdot a$ ($m$ times) & Definition of exponentiation \\
    $b^m = b \cdot b \cdot b\ldots \cdot b$ ($m$ times) & Definition of exponentiation \\
    $x^2 ~\rm{mod}~ q \iff (x ~\rm{mod}~ q)(x ~\rm{mod}~ q) ~\rm{mod}~ q$ & Law of modular multiplication for any arbitrary $x$. \\
    $a^m ~\rm{mod}~ n = (a ~\rm{mod}~ n)(a ~\rm{mod}~ n)\ldots (a ~\rm{mod}~ n) ~\rm{mod}~ n $ (m times) & Application of multiplication and exponentiation definitions \\
    $b^m ~\rm{mod}~ n = (b ~\rm{mod}~ n)(b ~\rm{mod}~ n)\ldots (b ~\rm{mod}~ n) ~\rm{mod}~ n $ (m times) & Application of multiplication and exponentiation definitions \\
    $b^m ~\rm{mod}~ n = (a ~\rm{mod}~ n)\ldots (a ~\rm{mod}~ n) ~\rm{mod}~ n $ (m times) & Since $a \equiv b ~\rm{mod}~ n$ \\
    $a^m \equiv b^m ~\rm{mod}~ n$ & Since we have established that both terms are equivalent to the multiplication expansion of $a^m ~\rm{mod}~ n$
\end{tabular}
\end{center}

Therefore, $a \equiv b ~\rm{mod}~ n \rightarrow a^m \equiv b^m ~\rm{mod}~ n$.

\newpage
\section{} % 5

\subsection{Provide the full statement of at least one theorem, as it’s stated on the presenter’s slides.}
Every $C^*$-algebra looks like a closed $*$-subalgebra of $B(H)$ for some Hilbert space $H$.

\subsection{In the theorem statement above, what type of a proposition is it? Is it conditional? Does it have suppressed quantifiers? If so, rewrite the theorem statement showing these.}
It is a universal statement because the theorem states that ``Every $C^*$-algebra...'' an example of this is:

\begin{center}
  $c(x) = \{f : x \rightarrow \mathbb{C} : $f is continuous$\}$ where x is compact, Hausdorff...
\end{center}

I suppose that the theorem could be conditional because it has to be a $C^*$-algebra.

We can write this as
\begin{center}
  $\forall C^* algebra~\exists H ~s.t. ~C^* algebra~ = ~closed~*$-subalgebra$~ \subseteq B(H)$.
\end{center}

\subsection{Write down a math term/word/notation that seems familiar to you (look out for words like injective, surjective, or notation like function notation).}
Let $(A(\mathbb{D}), \alpha)$ be a dynamical system. Then there exists a Mobius transformation of the form

\begin{center}
  $\tau(w) = \lambda (\frac{w-\mu}{1-\mu\bar w})$
\end{center}
such that $\alpha_1(f)(w) = f(\tau(w))$.

Within this I understand what it means for a system to be dynamical and for a transformation between spaces to exist.

\subsection{What do you understand this word/notation to mean?}
I understand this exerpt to mean there is some transform called a Mobius transform that exists as a function that accepts an argument $w$ and uses it to convert into something that is possibly called Mobius space.

\subsection{Using a resource (wikipedia is fine), provide a formal definition of the word or symbol as it was intended in the presentation.}
Here I look into the definition of a Mobius transformation.

In geometry and complex analysis, a Möbius transformation of the complex plane is a rational function of the form

\begin{center}
  $f(z) = \frac{az+b}{cz+d}$
\end{center}
of one complex variable $z$; here the coefficients $a,b,c,d$ are complex numbers satisfying $ab-bc \neq 0$.

"A Mobius transformation is always a bijective holomorphic function from the Riemann sphere to the Riemann sphere."

"These transformations preserve angles, map every straight line to a line or circle, and map every circle to a line or circle."

All information is probided from (Wikipedia, 2020). These exerpts provide the definition of a Mobius transformation as well as an invarient within the transformation.

\subsection{Write down a math term/word/notation that was totally unfamiliar to you.}
Irrational Rotation.

\subsection{Using a resource (wikipedia is fine), provide a formal definition of the word or symbol as it was intended in the presentation.}
In the mathematical theory of dynamicval systems, an irrational rotation is a map where $\theta$ is an irrational number. Under the identification of a circle with $R/z$ or with the interval $[0,1]$ with the boundary points glued together, this map becomes a rotation of a circle by a proportion $\theta$ of a full revolution.

\subsection{Ask a question that the presentation raised for you that you would like to know more about.}
Why isn't any orbit periodic? Within the same slide, Dr. Hamidi makes a statement "No orbit is periodic." What is a non-periodic orbit?


\newpage
\section{} % 6
POD 16: Below you will encrypt and decrypt messages using the RSA system.

\subsection{Provide the full statement of at least one theorem, as it’s stated on the presenter’s slides.}
Every $C^*$-algebra looks like a closed $*$-subalgebra of $B(H)$ for some Hilbert space $H$.

\subsection{In the theorem statement above, what type of a proposition is it? Is it conditional? Does it have suppressed quantifiers? If so, rewrite the theorem statement showing these.}
It is a universal statement because the theorem states that ``Every $C^*$-algebra...'' an example of this is:

\begin{center}
  $c(x) = \{f : x \rightarrow \mathbb{C} : $f is continuous$\}$ where x is compact, Hausdorff...
\end{center}

I suppose that the theorem could be conditional because it has to be a $C^*$-algebra.

We can write this as
\begin{center}
  $\forall C^* algebra~\exists H ~s.t. ~C^* algebra~ = ~closed~*$-subalgebra$~ \subseteq B(H)$.
\end{center}

\subsection{Write down a math term/word/notation that seems familiar to you (look out for words like injective, surjective, or notation like function notation).}
Let $(A(\mathbb{D}), \alpha)$ be a dynamical system. Then there exists a Mobius transformation of the form

\begin{center}
  $\tau(w) = \lambda (\frac{w-\mu}{1-\mu\bar w})$
\end{center}
such that $\alpha_1(f)(w) = f(\tau(w))$.

Within this I understand what it means for a system to be dynamical and for a transformation between spaces to exist.

\subsection{What do you understand this word/notation to mean?}
I understand this exerpt to mean there is some transform called a Mobius transform that exists as a function that accepts an argument $w$ and uses it to convert into something that is possibly called Mobius space.

\subsection{Using a resource (wikipedia is fine), provide a formal definition of the word or symbol as it was intended in the presentation.}
Here I look into the definition of a Mobius transformation.

In geometry and complex analysis, a Möbius transformation of the complex plane is a rational function of the form

\begin{center}
  $f(z) = \frac{az+b}{cz+d}$
\end{center}
of one complex variable $z$; here the coefficients $a,b,c,d$ are complex numbers satisfying $ab-bc \neq 0$.

"A Mobius transformation is always a bijective holomorphic function from the Riemann sphere to the Riemann sphere."

"These transformations preserve angles, map every straight line to a line or circle, and map every circle to a line or circle."

All information is probided from (Wikipedia, 2020). These exerpts provide the definition of a Mobius transformation as well as an invarient within the transformation.

\subsection{Write down a math term/word/notation that was totally unfamiliar to you.}
Irrational Rotation.

\subsection{Using a resource (wikipedia is fine), provide a formal definition of the word or symbol as it was intended in the presentation.}
In the mathematical theory of dynamicval systems, an irrational rotation is a map where $\theta$ is an irrational number. Under the identification of a circle with $R/z$ or with the interval $[0,1]$ with the boundary points glued together, this map becomes a rotation of a circle by a proportion $\theta$ of a full revolution.

\subsection{Ask a question that the presentation raised for you that you would like to know more about.}
Why isn't any orbit periodic? Within the same slide, Dr. Hamidi makes a statement "No orbit is periodic." What is a non-periodic orbit?


\end{document}
